
\begin{abstract}
Modern GPU-centric datacenter systems run diverse workloads, such as ML inference, training, and vector databases, that increasingly demand adaptive, workload-specific memory management, scheduling decisions, and fine-grained observability. However, today's GPU stacks embed these critical policies within isolated frameworks, closed-source vendor libraries, one-size-fit-all Linux kernel modules and device side code, making it prohibitively difficult to implement and experiment with new adaptive policies. While prior systems research has successfully extend CPU-side policies (e.g., Linux schedulers, page caches, congestion control), adapting these ideas directly to GPU-centric environments is challenging. First, retrofitting mechanism-policy separation into existing monolithic GPU drivers is non-trivial, as resource management logic is tightly coupled with low-level hardware interactions and lacks extension points. Second, CPU-side programmability alone is insufficient, as it cannot directly customize decisions within GPU kernels; likewise, GPU-side programmability alone is not enough, as it lacks coordination with host-side management and shared infrastructure. To overcome these limitations, we introduce \sys, a unified, safe, and efficient programmable policy runtime substrate that seamlessly spans both the CPU host and GPU devices. \sys refactors the monolithic GPU driver to decouple resource management mechanisms from dynamic policies, and offloads verified eBPF programs to run directly inside device side GPU kernels via a lightweight, sandboxed JIT runtime, enabling transparent, fine-grained control over memory offloading, kernel scheduling, and observability policies. Our evaluation across realistic workloads including LLM inference, GNN training and vector search demonstrates that \sys improve performance by up to 5x, improves tail latency by 2x, and adds negligible overhead. These results highlight the practical and essential nature of unified CPU--GPU extensibility for future GPU-centric systems.
\end{abstract}

\section{Introduction}

GPUs have become indispensable in modern datacenters, supporting an increasingly diverse and demanding set of workloads, from large-scale deep learning training and latency-sensitive inference to irregular computations like graph analytics and graphics rendering~\cite{kousha2019designing,shen2018cudaadvisor,zhou2021measurement,stocks2024breaking,das2023large,fedeli2022pushing,liu2021closing,guan2024amanda,kwon2023efficient,pankratz2021vulkan,sellers2016vulkan}. Effective management of GPU-centric systems requires fine-grained, adaptive policies spanning three critical domains: memory management and offloading, GPU kernel scheduling~\cite{hong2017gpu,kato2011timegraph}, and comprehensive observability of performance-critical events. However, despite their central importance, policies governing these domains are deeply embedded within rigid heuristic-based frameworks, proprietary vendor libraries, and opaque hardware and kernel-module interfaces. As workloads evolve, the inflexibility of these embedded policies leads directly to unpredictable performance, excessive latency, and costly downtime due to disruptive policy adjustments and system restarts.

% needs rewrite and fix to avoid duplication

In recent years, the broader systems community has demonstrated significant benefits from externalizing and customizing OS policies through frameworks like extended Berkeley Packet Filter (eBPF)~\cite{eBPFWindows}. Examples such as sched\_ext for Linux scheduling, cache\_ext for page-cache eviction, and struct\_ops for TCP congestion control have all illustrated substantial performance and adaptability gains when policies can be dynamically loaded, safely verified, and executed directly within kernel contexts~\cite{gregg16,gregg2021computing,Bachl21,jia2023practical,Zhong22,yang2023lambda,jia2023programmable,ghigoff2021bmc,chen2025etran,xu2025state,benson2024netedit,zhou2024dint,wu2025vep,bpftime,mao2024merlin,jin2024beebox,wang2024seak,zhong2022xrp}. However, while these successes illustrate the potential of policy programmability on CPUs, directly extending these ideas to GPU systems poses significant new challenges. CPU-side programmability alone is insufficient, as critical GPU management decisions, such as fine-grained control over GPU kernel admission, prefetching and eviction policies within page fault handlers, and precise on-device instrumentation, are fundamentally embedded in GPU device execution contexts inaccessible from userspace, tightly coupled with low-level interactions and lacks extension points. Conversely, GPU-side programmability alone is not enough either, as isolated GPU-only policies lack visibility into broader system state, host-side control, and coordination essential for coherent end-to-end policy decisions.

Current GPU drivers enforce rigid, ``one-size-fits-all'' policies (e.g., LRU eviction) that fail to adapt to the conflicting needs of modern workloads: a latency-sensitive LLM inference query requires aggressive prefetching of weights, while a random-access GNN training job requires thrashing-resistant caching. Emerging user-space runtimes (e.g., vLLM, Pie) attempt to bypass this by implementing custom management logic in user land. However, these approaches are fundamentally limited by a lack of privilege: they cannot intervene at the granularity of hardware interrupts (e.g., handling a page fault in microseconds) or directly manipulate physical memory mappings. Consequently, they miss critical opportunities for overlapping I/O with computation and cannot enforce true preemption in multi-tenant environments. Thus, extending the driver itself to expose these privileged mechanisms is the only path to transparent, fine-grained optimization.

While prior research has demonstrated substantial improvements from adaptive page placement, prefetching, and eviction policies in non-GPU contexts, such as key-value stores, CDNs, and OS-level page caches (as demonstrated by cache\_ext), adapting these sophisticated eviction algorithms directly into GPU is exceedingly difficult due to the complexity and opacity of modifying proprietary GPU drivers and the lack of programmable hooks at critical decision points. Similar limitations apply to GPU scheduling: existing GPU admission and resource allocation policies remain rigid, coarse-grained, and buried deep within vendor runtime logic. Moreover, current GPU observability tools like NVIDIA CUPTI and NVBit~\cite{cupti,villa2019nvbit} provide detailed monitoring data but lack programmability and actionable integration into policy decisions. Thus, while CPU-side policy programmability has proven transformative in purely CPU-oriented contexts, extending this programmability effectively to GPU systems demands a fundamentally new, unified CPU--GPU approach.


To address these critical limitations, we introduce \sys, a novel, unified framework that provides safe, efficient, and coherent programmability across both CPU and GPU domains. At its core, \sys extends the Linux \sysdriver to safely expose programmable, eBPF-based policy attach points directly at key decision-making sites, such as memory placement decisions, CUDA kernel scheduling logic, and in-kernel observability instrumentation. This architecture enables \textbf{transparent} optimizations that are impossible in user space, allowing policies to react to hardware events (like page faults) in microseconds. Additionally, \sys introduces a lightweight, sandboxed runtime that offloads and executes verified eBPF programs directly inside GPU kernels through launch-time PTX-level injection. By carefully combining CPU-side and GPU-side execution of verified, bounded eBPF programs, \sys enables workload-specific and adaptive policy logic to run exactly at the points in the system where decisions naturally occur, preserving both correctness and performance isolation.

Designing and implementing \sys introduced multiple technical challenges. First, extending open-source GPU driver modules to expose programmable hooks without compromising stability required careful engineering and minimal modifications to existing trusted paths. Second, safely offloading eBPF logic to GPUs posed novel verification and sandboxing constraints, ensuring no interference with user kernels, bounded resource usage, and minimal overhead. Third, developing a unified policy-state abstraction---shared eBPF maps accessible across CPU and GPU boundaries---was critical to ensure coherent state visibility and minimal synchronization overhead. \sys addresses these challenges through a carefully engineered combination of minimal kernel-module extensions, a GPU-side PTX-JIT policy runtime, and cross-domain eBPF map-based state sharing.

We evaluate \sys across realistic GPU workloads, including LLM inference, graph neural network training, vector search and mixed-priority tenant scenarios, running on modern GPU systems. Our results demonstrate the practical benefits of unified CPU--GPU programmability: compared to static heuristics, \sys-enabled adaptive policies; and observability instrumentation provides detailed actionable telemetry at reasonable overheads. Crucially, these improvements are achieved without disrupting workloads or requiring system restarts, highlighting \sys's practical impact.

\begin{itemize}[leftmargin=*,itemsep=0pt]
    \item We identify and characterize critical policy programmability gaps in GPU-centric stacks, specifically highlighting limitations in memory management, GPU scheduling, and observability due to embedded static heuristics, opaque vendor modules, and rigid GPU runtimes.

    \item We introduce \sys, the first system to offer unified, safe, efficient, and coherent programmability that spans CPU and GPU execution contexts. \sys extends the \sysdriver and provides on-device execution of verified eBPF programs directly within GPU kernels.

    \item We demonstrate through comprehensive evaluation that \sys's unified programmability significantly improves workload performance, resource utilization, and operational flexibility in realistic GPU-based datacenter scenarios, thereby opening the GPU stack to adaptive, workload-specific customization.
\end{itemize}

Ultimately, by bridging the divide between CPU and GPU programmability, \sys provides a critical foundation for flexible, efficient, and dynamically adaptive GPU-centric systems.


